% ==========================================================================
%  Chapter 7: Spatial Vector Algebra and Inertia
% ==========================================================================
\chapter{Spatial Vector Algebra and Inertia}
\label{ch:spatial_algebra}

\epigraph{%
  Why manage six tiny numbers loosely floating in space when you can forge them into an unbreakable, 6-dimensional spear?
}{}

\section{Featherstone's Spatial Algebra}

In Chapters 3 and 4, we unified translation and rotation into 6-dimensional Twists ($\twist$) and Wrenches ($\wrench$). While the geometry of the Screw Axis\index{Screw Axis} inherently links them, operating on these 6-dimensional vectors using standard 3D linear algebra creates a bloated, confusing mess of cross products and offsets.

To fix this, Roy Featherstone formalized \textbf{Spatial Vector Algebra}. Rather than writing Newton's equations as two separate 3D systems (one for $F = ma$ and one for $\tau = I\alpha + \omega \times I\omega$), Spatial Algebra unifies them into a single, breathtakingly elegant 6D format.

\begin{laymansbox}{The Spatial Paradigm}{featherstone_intuition}
  In 3D mechanics, linear motion ($v$) and angular motion ($\omega$) behave differently. In Spatial Algebra, they are fused. A "Spatial Vector" does not simply have a direction and a magnitude; it fundamentally exists as a line floating in the universe with a pitch. When math is done in this space, the separation between "turning" and "pushing" disappears completely.
\end{laymansbox}

\section{The Spatial Cross Product}

In classical 3D mechanics, the cross product $a \times b$ finds a perpendicular vector. In 6D Spatial Algebra, because Twists ($\twist$) and Wrenches ($\wrench$) transform differently under coordinate frame shifts, there are actually \textit{two} spatial cross products: one for motion vectors (like velocity) and one for force vectors.

The spatial cross product for motion ($\times$) is defined as a $6 \times 6$ matrix operator. If a frame has an angular velocity $\omega$ and linear velocity $v$, its Spatial Motion Cross operator $[\twist \times]$ is:

\begin{equation}
   [\twist \times] = \begin{bmatrix} [\omega] & 0 \\ [v] & [\omega] \end{bmatrix} \in \Reals^{6 \times 6}
\end{equation}
where $[\omega]$ and $[v]$ are the standard $3 \times 3$ skew-symmetric matrices.

The spatial force cross product ($\times^*$) is simply the negative transpose of the motion operator:
\begin{equation}
   [\twist \times^*] = -[\twist \times]^T = \begin{bmatrix} [\omega] & [v] \\ 0 & [\omega] \end{bmatrix}
\end{equation}

By keeping forces and motions in their respective dual spaces, calculating the Coriolis forces tearing at the joints of a swinging golf club requires only a single line of 6D matrix multiplication rather than pages of complex calculus.

\section{The Spatial Inertia Matrix}

Just as mass $m$ resists linear acceleration, and a $3 \times 3$ inertia tensor $I$ resists angular acceleration, standard 6D Spatial Algebra unifies them into a single $6 \times 6$ \textbf{Spatial Inertia Matrix} $I_s$.

For a rigid body with mass $m$, a center of mass located at vector $\bm{c}$, and a $3 \times 3$ rotational inertia tensor $I_c$ about that center of mass, the total Spatial Inertia Matrix $I_s$ about the origin of the frame is:

\begin{equation}
  I_s = \begin{bmatrix} I_c + m[\bm{c}][\bm{c}]^T & m[\bm{c}] \\ m[\bm{c}]^T & m \mathbf{1}_{3 \times 3} \end{bmatrix} \in \Reals^{6 \times 6}
\end{equation}

This single matrix fundamentally dictates how hard it is to alter the trajectory of a physical object crossing through any dimensional axis.

\section{Worked Example: Spatial Inertia of a Cylinder}

Let us build a computational example verifying these principles. In robot design, human body modeling (like the 15-DOF golfer), or physics engines (MuJoCo), we must frequently construct $6 \times 6$ inertia tensors from basic geometric primitive shapes. 

Here is the python code required to compute the Spatial Inertia Matrix for a solid cylinder (such as a robotic forearm or a humerus bone) of mass $m$, radius $r$, and length $L$, whose center of mass is shifted along the Z-axis by a distance $d$.

\begin{lstlisting}[language=Python, caption=Computing 6D Spatial Inertia for a Cylinder]
import numpy as np

def skew(v):
    # Converts a 3D vector into a 3x3 skew-symmetric matrix
    return np.array([
        [ 0,    -v[2],  v[1]],
        [ v[2],  0,    -v[0]],
        [-v[1],  v[0],  0   ]
    ])

def cylinder_spatial_inertia(m, r, L, d):
    # 1. Rotational inertia tensor at the Center of Mass (CoM)
    I_xx = (1/12) * m * (3*r**2 + L**2)
    I_yy = I_xx
    I_zz = (1/2) * m * r**2
    I_c = np.diag([I_xx, I_yy, I_zz])
    
    # 2. Center of Mass offset vector
    c = np.array([0.0, 0.0, d])
    c_skew = skew(c)
    
    # 3. Assemble the 6x6 Spatial Inertia Matrix
    I_s = np.zeros((6, 6))
    
    # Top Left: I_c + m * [c][c]^T (Parallel Axis Theorem)
    I_s[0:3, 0:3] = I_c + m * (c_skew @ np.transpose(c_skew))
    
    # Top Right: m * [c]
    I_s[0:3, 3:6] = m * c_skew
    
    # Bottom Left: m * [c]^T
    I_s[3:6, 0:3] = m * np.transpose(c_skew)
    
    # Bottom Right: m * Identity
    I_s[3:6, 3:6] = m * np.eye(3)
    
    return I_s

# Example: A 2kg forearm cylinder, 5cm radius, 30cm long, offset 15cm down
I_spatial = cylinder_spatial_inertia(m=2.0, r=0.05, L=0.30, d=0.15)
print("6x6 Spatial Inertia Matrix:\n", np.round(I_spatial, 4))
\end{lstlisting}

When you execute this code, the Parallel Axis Theorem is automatically handled natively inside the $6 \times 6$ matrix. The resulting matrix perfectly defines how hard that cylindrical arm is to accelerate, twist, wrench, or twist-push simultaneously around the joint.
